\section{Cognitive Loop}

The cognitive loop is the central operational cycle of the
MEM8-Geometric-Architecture. It defines how semantic information is perceived,
encoded, evaluated, routed, stabilized, and ultimately used to produce
behavioral output. Each stage corresponds to a geometric transformation on the
semantic manifold, ensuring continuity, stability, and interpretability.

Let $\mathbf{x}_t$ denote the raw input at time $t$, and let
$\mathbf{z}_t$ denote the fused semantic representation after encoding,
salience scoring, and stability filtering. The cognitive loop consists of six
stages.

\subsection{1. Perception}

Raw input is acquired from one or more modalities:
\[
\mathbf{x}_t^{(text)},\quad
\mathbf{x}_t^{(image)},\quad
\mathbf{x}_t^{(audio)}
\]

Each modality is mapped into a preliminary embedding space using
modality-specific encoders. This produces:
\[
\mathbf{t}_t,\quad \mathbf{i}_t,\quad \mathbf{a}_t \in \mathbb{R}^d
\]

\subsection{2. Encoding}

Encoding transforms raw modality embeddings into a geometric memory substrate.
Wave memory encoding is applied to each modality:
\[
\mathbf{X}_{wave} = \mathcal{F}(\mathbf{x})
\]

The fused semantic vector is then computed:
\[
\mathbf{z}_t = W_t \mathbf{t}_t + W_i \mathbf{i}_t + W_a \mathbf{a}_t
\]

\subsection{3. Salience Scoring}

Salience quantifies the behavioral relevance of the current semantic state.
Novelty is measured via Mahalanobis distance:
\[
N(\mathbf{z}_t) =
\sqrt{(\mathbf{z}_t - \bm{\mu})^\top \Sigma^{-1}
(\mathbf{z}_t - \bm{\mu})}
\]

Temporal change is measured as drift:
\[
T(\mathbf{z}_t) = \|\mathbf{z}_t - \mathbf{z}_{t-1}\|
\]

Combined salience:
\[
S(\mathbf{z}_t) =
\alpha N(\mathbf{z}_t) + (1 - \alpha) T(\mathbf{z}_t)
\]

\subsection{4. Stability Filtering}

Stability filtering smooths the semantic trajectory using a Kalman filter.

Prediction:
\[
\mathbf{z}_{t|t-1} = A \mathbf{z}_{t-1} + B \mathbf{u}_t
\]

Update:
\[
\mathbf{z}_{t|t} =
\mathbf{z}_{t|t-1} +
K_t (\mathbf{y}_t - \mathbf{z}_{t|t-1})
\]

This ensures geometric continuity and prevents collapse-risk escalation.

\subsection{5. Routing}

Routing selects the geometric direction most aligned with the stabilized
semantic state:
\[
i^\* = \arg\max_i\,
\cos(\mathbf{z}_{t|t}, \mathbf{r}_i)
\]

Each route vector $\mathbf{r}_i$ corresponds to a distinct cognitive operation.

\subsection{6. Action Selection}

The selected route determines the downstream operator applied to the semantic
state. This may include:

\begin{itemize}
\item memory update,
\item contextual expansion,
\item modality reweighting,
\item stability reinforcement,
\item output generation.
\end{itemize}

Action selection produces the final output:
\[
\mathbf{a}_t = f_{i^\*}(\mathbf{z}_{t|t})
\]

\subsection{Geometric Interpretation}

The cognitive loop is a closed geometric circuit. Each stage corresponds to a
transformation on the semantic manifold:

\begin{itemize}
\item perception: manifold entry,
\item encoding: geometric projection,
\item salience: curvature-sensitive relevance scoring,
\item stability: drift and gradient smoothing,
\item routing: directional alignment,
\item action: geodesic continuation.
\end{itemize}

This geometric formulation enables the MEM8-inspired cognitive architecture to
operate continuously, stably, and interpretably across diverse semantic inputs.
